and hence such that @@M0@@ is a @@M1@@-Module of finite type, and therefore of finite presentation by the noetherian hypotheses; the result moreover holds for every integer @@M2@@, since the @@M3@@ are coherent by Serre's theorem.

The proposition is therefore proved in the case where @@M4@@ is noetherian.

Suppose now that we have a Cartesian square:
@@M5@@
with @@M6@@ noetherian, @@M7@@; let @@M8@@ be a @@M9@@-Module of finite presentation, flat over @@M10@@, and let @@M11@@. By the preceding argument the proposition holds for @@M12@@. By EGA III 7.7.5 and 7.8.5, the base-change theorem holds under the preceding hypotheses; we immediately deduce that the theorem holds for @@M13@@.

To complete the proof, it therefore suffices to show that the general case (with @@M14@@ affine) can be obtained from the noetherian case by base change.

Since @@M15@@ is affine, @@M16@@ is a filtered projective limit of noetherian affine schemes (if @@M17@@, we write @@M18@@ as the inductive limit of its subalgebras of finite type over @@M19@@). Thus let @@M20@@. By EGA IV 8.5.2 and 8.5.5, there exists an index @@M21@@ and a @@M22@@-Module locally free of finite type @@M23@@ such that @@M24@@ (denoting by @@M25@@ the morphism @@M26@@, and by @@M27@@ the transition morphisms @@M28@@).

If we put, for @@M29@@, @@M30@@, and @@M31@@, then for every @@M32@@, we have @@M33@@, and @@M34@@. Consequently, there exists an index @@M35@@
